\begin{table}[H]
\centering
\caption{Nama Daerah Perbandingan Cluster Statis vs IPM Rata-rata (UNDP)} \label{tab:statis_ipm_avg_name}
\begin{adjustbox}{max width=\textwidth}
\begin{tabular}{|p{4cm}|p{10cm}|}
\hline
\textbf{Kategori} & \textbf{Daftar Wilayah} \\
\hline
Kl.1 \& Sangat Tinggi & - \\
\hline
Kl.1 \& Tinggi & Kep. Seribu, Ciamis, Kota Banjar \\
\hline
Kl.1 \& Sedang & Cianjur, Garut, Tasikmalaya, Majalengka, Indramayu, Pangandaran, Pandeglang, Lebak, Serang \\
\hline
Kl.1 \& Rendah & - \\
\hline
Kl.2 \& Sangat Tinggi & - \\
\hline
Kl.2 \& Tinggi & Bogor, Bandung, Kuningan, Cirebon, Sumedang, Subang, Purwakarta, Karawang, Kota Cirebon, Kota Tasikmalaya, Tangerang, Kota Cilegon, Kota Serang \\
\hline
Kl.2 \& Sedang & Sukabumi, Bandung Barat \\
\hline
Kl.2 \& Rendah & - \\
\hline
Kl.3 \& Sangat Tinggi & Kota Jakarta Selatan, Kota Jakarta Timur, Kota Jakarta Pusat, Kota Jakarta Barat, Kota Jakarta Utara, Kota Bandung, Kota Bekasi, Kota Depok, Kota Tangerang Selatan \\
\hline
Kl.3 \& Tinggi & Bekasi, Kota Bogor, Kota Sukabumi, Kota Cimahi, Kota Tangerang \\
\hline
Kl.3 \& Sedang & - \\
\hline
Kl.3 \& Rendah & - \\
\hline
\end{tabular}
\end{adjustbox}
\end{table}